\documentclass[12pt,a4paper]{article}
\usepackage[T1]{fontenc}
\usepackage[utf8]{inputenc}
\usepackage[ngerman,latin]{babel}
\usepackage{amsmath,amssymb,amsfonts}
\usepackage{geometry}
\usepackage{fancyhdr}
\usepackage{titlesec}
\usepackage{array}
\usepackage{booktabs}
\usepackage{etoolbox}
\usepackage{lettrine}

\geometry{
  paperwidth=170mm,
  paperheight=250mm,
  top=2.5cm,
  bottom=2.5cm,
  inner=2.5cm,
  outer=2cm,
  headheight=15pt
}

\tolerance=1000
\emergencystretch=3em
\hyphenpenalty=50

\title{Carl Friedrich Gauß Werke --- Band 5, Teil 14}
\author{Carl Friedrich Gauß}
\date{}

% Custom page style for Bemerkungen pages
\fancypagestyle{bemerkungen}{
  \fancyhf{}
  \fancyhead[C]{\textsc{Bemerkungen.}}
  \fancyfoot[RO,LE]{\thepage}
  \renewcommand{\headrulewidth}{0pt}
  \renewcommand{\footrulewidth}{0pt}
}

% Custom page style for Inhalt pages
\fancypagestyle{inhalt}{
  \fancyhf{}
  \fancyhead[C]{\textsc{Inhalt.}}
  \fancyfoot[RO,LE]{\thepage}
  \renewcommand{\headrulewidth}{0pt}
  \renewcommand{\footrulewidth}{0pt}
}

% Custom arc tang operator
\DeclareMathOperator{\arctang}{arc\,tang}

% Custom command for Inhalt entries
\newcommand{\inhaltentry}[3]{%
  \noindent\parbox[t]{0.72\textwidth}{\raggedright #1}%
  \hfill\parbox[t]{0.26\textwidth}{\raggedleft #2\quad---\quad#3}\par
  \vspace{0.3em}
}

% Custom command for Inhalt section headers
\newcommand{\inhaltsection}[1]{%
  \vspace{0.5em}
  \noindent\textit{#1}\par
  \vspace{0.3em}
}

\begin{document}

% ============================================
% Title page
% ============================================
\maketitle
\thispagestyle{empty}

\vfill

\begin{center}
\textsc{Carl Friedrich Gauß Werke.}\\[1em]
\textsc{Fünfter Band.}\\[1em]
\textsc{Mathematische Physik.}\\[2em]
{\small Herausgegeben von der Königlichen Gesellschaft der Wissenschaften zu Göttingen.}\\[2em]
{\large Göttingen.}\\
{\small Dieterichsche Universitäts-Buchdruckerei.}\\
{\small W.~Fr.~Kaestner.}\\
{\small 1867.}
\end{center}

\vfill

\newpage

% ============================================
% Page 639 (PDF 651) -- BEMERKUNGEN
% ============================================
\thispagestyle{bemerkungen}
\setcounter{page}{639}

\begin{center}
\textsc{Bemerkungen.}
\end{center}

\vspace{0.5em}

Dieser Satz gibt durch wiederholte Anwendung auch den Beweis von \textsc{Ampères} Fundamentalsatz in der allgemeinen Form, dass unmittelbar die Potentialfunction für die Wechselwirkung zwischen den auf bestimmte Weise mit magnetischem Fluidum belegten Flächen zurückgeführt wird auf die Potentialfunction für die Wechselwirkung zwischen galvanischen Strömen, die nach Lage und Intensität durch jene Flächen und die Magnetisirung bestimmt sind.

\vspace{0.5em}

Dem in Nr.~9. aufgestellten Beweise für die Gleichheit der Werthe der verschiedenen Ausdrücke für die Potentialfunction~$V$ kann man eine symmetrische Form geben, wenn man die Function~$R$
\begin{multline*}
= xx\arctang\frac{yz}{xr} + yy\arctang\frac{zx}{yr} + zz\arctang\frac{xy}{zr}
\\
{}+ 2yzi\arctang\frac{xi}{r} + 2zxi\arctang\frac{yi}{r} + 2xyi\arctang\frac{zi}{r}
\end{multline*}
worin $i$ statt $\sqrt{-1}$ gesetzt ist, einführt und berücksichtigt, dass die Gleichungen
\begin{align*}
\frac{d\,d\,R}{d\,x^2} &= 2\arctang\frac{yz}{xr}, &
\frac{d\,d\,R}{d\,y\,d\,z} &= 2i\arctang\frac{xi}{r}, \\
\frac{d\,d\,R}{d\,y^2} &= 2\arctang\frac{zx}{yr}, &
\frac{d\,d\,R}{d\,z\,d\,x} &= 2i\arctang\frac{yi}{r}, \\
\frac{d\,d\,R}{d\,z^2} &= 2\arctang\frac{xy}{zr}, &
\frac{d\,d\,R}{d\,x\,d\,y} &= 2i\arctang\frac{zi}{r}
\end{align*}
\[
\frac{d\,d\,R}{d\,x^2} + \frac{d\,d\,R}{d\,y^2} + \frac{d\,d\,R}{d\,z^2} = (2m+1)\pi,
\qquad
\frac{d^3 R}{d\,x\,d\,y\,d\,z} = -2\frac{1}{r}
\]
Statt haben. Durch die Derivirten der Function $R(x,y,z)$ können in endlicher Form auch die Potentialfunctionen für die magnetische Wirkung solcher galvanischer Ströme dargestellt werden, deren Leiter aus geradlinigen den Axen eines rechtwinkligen Coordinatensystems parallelen linearen Theilen bestehen.

\vspace{0.5em}

Die von \textsc{Gauss} bei der Bestimmung einer Abtheilung eines getheilten Maassstabes in Theilen des ganzen Maassstabes angewandte Anordnung der Längen-Comparirungen danken wir der Aufzeichnung, die sich der Herr Geh.~Hofrath \textsc{Weber} im Jahre 1839 oder 1840 gemacht hat.

\vspace{0.5em}

Die über die Beugungserscheinungen angestellten theoretischen Untersuchungen sind wahrscheinlich durch das von \textsc{F.~M.~Schwerd} im Jahre 1835 herausgegebene diesen Gegenstand betreffende Werk veranlasst. Die beiden für die Wirkung eines leuchtenden Punkts~$P$ auf einen Punkt~$p$ aufgestellten allgemeinen Formeln sind nicht identisch; die allgemeine Verwandlung solcher Flächen-Integrale, deren Elemente von der Lage der durch einen Punkt des zugehörigen Flächentheilchens~$ds$ und durch die Punkte~$P$ und~$p$ gehenden Ebene nicht abhängen, in solche Curven-Integrale, deren Elemente ebenfalls von der Lage der durch einen Punkt des zugehörigen Theilchens der Begrenzungslinie~$u$ und durch die Punkte~$P$ und~$p$ gehenden Ebene nicht abhängen, deren Differentiale aber eine Änderung allein des Winkels~$\vartheta$ bedeuten, welchen jene Ebene mit einer durch~$P$ und~$p$ gelegten festen Ebene einschliesst, ergibt sich aus der Gleichung

% ============================================
% Page 640 (PDF 652) -- BEMERKUNGEN (continued)
% ============================================
\newpage
\thispagestyle{bemerkungen}
\setcounter{page}{640}

\begin{align*}
\int Q\,d\vartheta &= \int\frac{h\,Q\sin v\,du}{rR\sin w} \\
&= \int\!\frac{h}{rR}\,\frac{1}{2\sin\!\frac{1}{2}w^2}\,\frac{dQ}{d(R-r)}\,\frac{d(R+r)}{d\rho}\,ds \\
&\quad - \int\!\frac{h}{rR}\,\frac{1}{2\cos\!\frac{1}{2}w^2}\,\frac{dQ}{d(R+r)}\,\frac{d(R-r)}{d\rho}\,ds
\end{align*}
die einen speciellen Fall des in der vorhergehenden Bemerkung erwähnten Satzes bildet, wenn nemlich $Q$ eine von $R$ und $r$ allein abhängige mit ihren nach $R+r$ und $R-r$ genommenen partiellen Derivirten für die Punkte der Fläche~$s$ stetig veränderliche Grösse bedeutet, ferner~$\varpi$ genauer als im Text dahin bestimmt ist, dass es den Winkel zwischen $R$ und~$r$ bezeichnet, den die mit den Richtungen des fortschreitenden Lichtstrahls übereinstimmend angenommenen positiven Richtungen jener Linien einschliessen.

\vspace{2em}

\hfill\textsc{Schering.}

\vfill

\begin{center}
\rule{0.3\textwidth}{0.4pt}
\end{center}

\newpage

% ============================================
% Page 641 (PDF 653) -- INHALT (first page)
% ============================================
\thispagestyle{inhalt}
\setcounter{page}{641}

\begin{center}
\textsc{Inhalt.}
\end{center}

\vspace{1em}

\begin{center}
\textbf{GAUSS WERKE BAND V. MATHEMATISCHE PHYSIK.}
\end{center}

\vspace{1em}

\inhaltsection{Abhandlungen.}

\inhaltentry{Theoria attractionis corporum sphaeroidicorum ellipticorum homogeneorum}{1813 März\ldots}{Seite~1}

\inhaltentry{Über ein neues allgemeines Grundgesetz der Mechanik}{1829}{---~23}

\inhaltentry{Principia generalia theoriae figurae fluidorum in statu aequilibrii}{1829 Sept.\ldots}{---~29}

\inhaltentry{Intensitas vis magneticae terrestris ad mensuram absolutam revocata}{1832 Dec.\ldots}{---~79}

\inhaltentry{Allgemeine Theorie des Erdmagnetismus}{1838}{---~119}

\inhaltentry{Allgemeine Lehrsätze in Beziehung auf die im verkehrten Verhältnisse des Quadrats der Entfernung wirkenden Anziehungs- und Abstossungs-Kräfte}{1839}{---~195}

\inhaltentry{Dioptrische Untersuchungen}{1840 Dec.\ldots}{---~243}

\inhaltsection{Anzeigen eigner Abhandlungen.}

\inhaltentry{Theoria attractionis corporum sphaeroidicorum ellipticorum homogeneorum}{1813 April.\ldots}{---~279}

\inhaltentry{Principia generalia theoriae figurae fluidorum in statu aequilibrii}{1829 Oct.\ldots}{---~287}

\inhaltentry{Intensitas vis magneticae terrestris ad mensuram absolutam revocata}{1832 Dec.\ldots}{---~293}

\inhaltentry{Allgemeine Lehrsätze in Beziehung auf die im verkehrten Verhältnisse u.~s.~w.}{1840 März.\ldots}{---~305}

\inhaltentry{Dioptrische Untersuchungen}{1841 Jan.\ldots}{---~309}

\inhaltsection{Verschiedene Aufsätze über Magnetismus.}

\inhaltentry{Erdmagnetismus und Magnetometer}{1836}{---~315}

\inhaltentry{Einleitung für die Zeitschrift: Resultate u.~s.~w.}{1836}{---~345}

\inhaltentry{Ein neues Hülfsmittel für die magnetischen Beobachtungen}{1837 Oct.\ldots}{---~352}

\noindent\parbox[t]{0.72\textwidth}{\raggedright
Über ein neues, zunächst zur unmittelbaren Beobachtung der Veränderungen in der Intensität des horizontalen Theils des Erdmagnetismus bestimmten Instruments
}\hfill\parbox[t]{0.26\textwidth}{\raggedleft 1837\quad---\quad357}\par
\vspace{0.3em}

\inhaltentry{Anleitung zur Bestimmung der Schwingungsdauer einer Magnetnadel}{1837}{---~374}

\inhaltentry{Über ein Mittel die Beobachtung von Ablenkungen zu erleichtern}{1839}{---~395}

\inhaltentry{Zur Bestimmung der Constanten des Bifilarmagnetometers}{1840}{---~404}

\noindent\parbox[t]{0.72\textwidth}{\raggedright
Vorschriften zur Bestimmung der magnetischen Wirkung, welche ein Magnetstab in der Ferne ausübt
}\hfill\parbox[t]{0.26\textwidth}{\raggedleft 1840\quad---\quad427}\par
\vspace{0.3em}

\vfill

\begin{flushright}
V.
\end{flushright}

% ============================================
% Page 642 (PDF 654) -- INHALT (continued)
% ============================================
\newpage
\thispagestyle{inhalt}
\setcounter{page}{642}

\inhaltentry{Über die Anwendung des Magnetometers zur Bestimmung der absoluten Declination}{1841}{Seite~436}

\inhaltentry{Beobachtungen der magnetischen Inclination in Göttingen}{1841}{---~444}

\inhaltsection{Aufsätze über verschiedene Gegenstände der mathematischen Physik.}

\noindent\parbox[t]{0.72\textwidth}{\raggedright
Fundamentalgleichungen für die Bewegung schwerer Körper auf der rotirenden Erde
}\hfill\parbox[t]{0.26\textwidth}{\raggedleft 1803\quad---\quad495}\par
\vspace{0.3em}

\noindent\parbox[t]{0.72\textwidth}{\raggedright
Über die achromatischen Doppelobjective besonders in Rücksicht der vollkommenern Aufhebung der Farbenzerstreuung
}\hfill\parbox[t]{0.26\textwidth}{\raggedleft 1817 Dec.\ldots\quad---\quad504}\par
\vspace{0.3em}

\inhaltentry{Brief an \textsc{Brandes} über denselben Gegenstand}{1831}{---~509}

\inhaltentry{Berichtigung der Stellung der Schneiden einer Wage}{1837 März.\ldots}{---~511}

\inhaltsection{Physicalische Beobachtungen.}

\inhaltentry{Nordlicht am 7. Januar 1831}{1831 Febr.\ldots}{---~517}

\inhaltentry{Magnetisches Observatorium in Göttingen}{1834 Aug.\ldots}{---~519}

\inhaltentry{Beobachtungen der magnetischen Variation in Göttingen und Leipzig}{1834 Oct.\ldots}{---~525}

\inhaltentry{Magnetisches Observatorium in Göttingen}{1835 März.\ldots}{---~528}

\inhaltentry{Beobachtungen der magnetischen Variation in Copenhagen und Mailand}{1835 März.\ldots}{---~537}

\inhaltentry{Magnetisches Observatorium in Göttingen}{1836 Juni.\ldots}{---~540}

\inhaltentry{Das in den Beobachtungsterminen anzuwendende Verfahren}{1836}{---~541}

\noindent\parbox[t]{0.72\textwidth}{\raggedright
Auszug aus dreijährigen täglichen Beobachtungen der magnetischen Declination zu Göttingen
}\hfill\parbox[t]{0.26\textwidth}{\raggedleft 1836\quad---\quad556}\par
\vspace{0.3em}

\inhaltentry{Erläuterungen zu den Terminszeichnungen und den Beobachtungszahlen}{1836}{---~568}

\inhaltentry{Erläuterungen zu den Terminszeichnungen und den Beobachtungszahlen}{1837}{---~576}

\inhaltentry{Der magnetische Südpol der Erde}{1841}{---~580}

\inhaltsection{Anzeigen nicht eigner Schriften.}

\inhaltentry{\textsc{Benzenberg.} Über die \textsc{Dalton}sche Theorie}{1830 Dec.\ldots}{---~583}

\inhaltentry{\textsc{Fischer.} Künstliche Magnete}{1832 Sept.\ldots}{---~591}

\inhaltentry{Resultate aus den Beobachtungen des magnetischen Vereins}{1837 Juni.\ldots}{---~595}

\inhaltentry{\textsc{Gerling.} Einrichtung des mathematisch-physicalischen Instituts}{1842 April.\ldots}{---~596}

\inhaltsection{Nachlass.}

\inhaltentry{Zur Electrodynamik}{}{---~601}

\inhaltentry{Über Kugelfunctionen}{}{---~630}

\inhaltentry{Zum Gebrauch des Comparators}{}{---~632}

\noindent\parbox[t]{0.72\textwidth}{\raggedright
Allgemeine Formeln für die Wirkung eines leuchtenden Punkts~$P$ auf einen Punkt~$p$
}\hfill\parbox[t]{0.26\textwidth}{\raggedleft \quad---\quad635}\par
\vspace{0.3em}

\inhaltentry{Bemerkungen von \textsc{Schering}}{}{---~637}

\vspace{0.3em}

\inhaltentry{Steindrucktafel zur Theorie des Erdmagnetismus Seite~176.}{}{}

\vfill

\begin{center}
\rule{0.5\textwidth}{0.4pt}
\end{center}

\vspace{2em}

\begin{center}
\textbf{GÖTTINGEN,}\\[0.5em]
\textbf{DRUCK DER DIETERICHSCHEN UNIVERSITÄTS-BUCHDRUCKEREI.}\\[0.5em]
\textbf{W.~FR.~KAESTNER.}
\end{center}

\newpage

% ============================================
% Pages 643--646 (PDF 655--658) -- Blank pages
% ============================================
\thispagestyle{empty}
\setcounter{page}{643}
\mbox{}
\newpage
\thispagestyle{empty}
\setcounter{page}{644}
\mbox{}
\newpage
\thispagestyle{empty}
\setcounter{page}{645}
\mbox{}
\newpage
\thispagestyle{empty}
\setcounter{page}{646}
\mbox{}

% ============================================
% Page 647 (PDF 659) -- Library card / end matter
% ============================================
\newpage
\thispagestyle{empty}
\setcounter{page}{647}

\vfill

\begin{center}
\textit{[Universitätsbibliothek Berkeley --- Leihkarte]}
\end{center}

\begin{center}
\fbox{\parbox{0.6\textwidth}{
\centering
\textbf{UNIVERSITY OF CALIFORNIA LIBRARY}\\
\textbf{BERKELEY}\\[1em]
Return to desk from which borrowed.\\[0.5em]
This book is DUE on the last date stamped below.\\[1em]
\begin{tabular}{|c|c|c|}
\hline
M677 & ASTRONOMY LIBRARY & \\
\hline
APR 1 & '66 & \\
\hline
& & \\
\hline
& & \\
\hline
\end{tabular}
}}
\end{center}

\vfill

% ============================================
% Page 648 (PDF 660) -- Back cover / endpaper
% ============================================
\newpage
\thispagestyle{empty}
\setcounter{page}{648}

\vfill

\begin{center}
\textit{[Marmoriertes Vorsatzpapier / Einbandrückseite]}
\end{center}

\vfill

\end{document}
